\section{Conclusion}
Bayesian spectral decomposition of extended b-value prostate DWI explains why ADC is so hard to improve on for ROI-level tumor detection. The tumor--normal contrast this acquisition can see is concentrated in the restricted-diffusion and free-water compartments, which are also the compartments ADC weights most heavily, so ranking regions by ADC reproduces the ranking of the optimal spectral classifier.
The severe ill-conditioning of the inverse problem leaves the intermediate diffusivities indistinguishable even at the low noise of an endorectal-coil acquisition, and the Bayesian treatment is what makes that limit legible, separating the compartments the data determine from those fixed largely by the prior and the non-negativity constraint.
Its contribution is therefore upstream of the diagnostic decision rather than within it, marking where a spectral readout could carry information ADC does not. Grade-dependent spectral shifts and pixel-wise maps with per-voxel uncertainty indicate where that might be, and both require larger studies with histopathological correlation before they can be claimed.
